\section{Introduction}\label{sec:intro}

A consistent lesson has emerged from the recent long-horizon time-series forecasting (LTSF) literature: performance is driven less by architectural novelty than by a small set of design components. Single-layer linear models match or beat elaborate Transformer variants on standard benchmarks \citep{zeng2023dlinear}; \citet{li2023rlinear} show that a linear map equipped with reversible instance normalization and channel independence suffices to outperform far larger models; \citet{toner2024linear} prove that several popular normalization-augmented linear forecasters are functionally equivalent to unconstrained linear regression with closed-form solutions; and the ablations of \citet{nie2023patchtst} attribute the gains of PatchTST primarily to patching and channel handling rather than attention. Large-scale meta-studies reinforce the point: no architecture wins consistently across datasets \citep{brigato2026champions}, and uncontrolled design dimensions---normalization, splits, even the \texttt{drop\_last} flag---can reorder leaderboards on their own \citep{qiu2024tfb, position2025benchmark}. If components, not architectures, decide performance, then the open question is no longer \emph{which component helps} but \emph{under which conditions each component's assumptions hold}.

Among these components, instance normalization is arguably the most decisive. RevIN \citep{kim2022revin} removes per-window mean and variance before the backbone and restores them afterward, neutralizing distribution shift between training and test; successors learn future distribution coefficients \citep{fan2023dishts}, predict slice-level statistics \citep{liu2023san}, or handle seasonal non-stationarity in the frequency domain \citep{ye2024fan}. Crucially, every member of this family is \emph{endogenous}: the statistics restored at the horizon are estimated from past observations of the target series itself. Recent work has begun to document situations where this backfires \citep{role2026revin, noise2025signal, inner2025instance}, and \citet{noise2025signal} go furthest, profiling each dataset with endogenous statistics and pre-selecting a normalization by an explicit threshold rule on change-point risk. That rule, however, chooses \emph{within} the endogenous family---standard versus robust instance normalization---using the target's own past, its threshold is a free hyperparameter set by hand, and its authors report that the resulting selector fails systematically, closing with a call for a principled diagnostics-driven strategy. What is still missing is therefore twofold: a conditional, quantitative characterization of \emph{when} and \emph{by how much} instance normalization hurts, and a pre-training diagnostic that decides between endogenous and \emph{exogenous} normalization statistics---that is, one whose choice set includes not estimating the level from the window at all. \Cref{sec:related} surveys this landscape in detail.

The gap matters most where series levels are \emph{exogenously determined}. In energy forecasting---the setting of the GEFCom competitions \citep{hong2016gefcom}---the level of wind generation is set by installed capacity times the weather regime, observable in advance through numerical weather prediction (NWP) archives; cooling load responds to temperature through a cooling-degree-day relationship; solar output follows a clear-sky structure modulated by cloud cover. For such series, the window mean that instance normalization subtracts is not a nuisance statistic to be discarded: it carries signal that is recoverable from covariates available at forecast time. Removing it forces the model to re-extrapolate the level from a quantity---the lookback-window mean---that is stale precisely when the level is about to move, and the staleness compounds with the forecast horizon. The repair is not in dispute and is not new: condition the level on the covariates, by a first-stage regression, by a classical dynamic regression, or---the instantiation we use as an \emph{instrument} throughout this paper---by covariate-based \emph{conditional normalization} (CondNorm), which normalizes by a first-stage regression of the level on exogenous covariates rather than by window statistics. What is at issue is the layer that makes the repair necessary. Accordingly we do not propose CondNorm as a method, and \Cref{sec:empirical} reports that two cheaper covariate-conditioned alternatives match or beat it on the very series where the rule fires; that is consistent with our claim rather than a concession against it, because the claim is about the normalization layer and not about any one replacement for it.

We therefore ask three questions. \textbf{RQ1}: Under what conditions does instance normalization lose to global z-scoring or to conditional normalization, and can the crossover be characterized in closed form? \textbf{RQ2}: How large is the effect on real data, relative to the choice of backbone? \textbf{RQ3}: Can the winning normalization be predicted \emph{before training} from the data alone?

Our contributions map onto these questions:
\begin{itemize}
  \item \textbf{C1 (Boundary).} In a linear-Gaussian instance model we give closed-form multi-step risks for global $z$-scoring (\rawm), instance normalization (\inm), and covariate-conditional normalization (\cnm), and read the boundary off their structural asymmetry: \inm{} is immune to level variance and to covariate-orthogonal distribution shift---the property it was designed for---but pays the full within-horizon level displacement plus a window-noise floor, so its advantage falls as the exogenous level predictability $\lambda$ rises and, for fixed $\lambda$, as the horizon grows (\Cref{sec:theory}). A synthetic grid of 2{,}640 runs reproduces the predicted crossings to within $0.015$ (\Cref{sec:synthetic}).
  \item \textbf{C2 (Mechanism).} The effect is a property of the \emph{layer}, not of covariate access. Under information parity---identical past and future covariates supplied to every arm across five covariate-capable backbones---instance normalization is worse in mean MSE than plain global $z$-scoring on \emph{all five}, so no amount of downstream flexibility recovers a level the normalizer has already removed. Correspondingly, three different covariate-conditioned level models (a first-stage regression alone, a classical dynamic regression, and \cnm) beat the best instance-normalized arm in 11 of 12 comparisons on the exogenously driven group and in 0 of 12 on the standard benchmarks (\Cref{sec:empirical}).
  \item \textbf{C3 (Coordinate).} The Level Predictability Score $\lps$, an out-of-sample $R^2$ of window-mean levels regressed on covariates under chronological cross-validation, computable before training. Its role is to place a series on one side of the boundary---$\lps \geq \tau$ means \emph{do not normalize the level endogenously}---and the sign it predicts was correct on all eight datasets of a study whose per-dataset predictions and protocol were committed to version control before any grid run, unchanged under MAE and MASE (\Cref{sec:lps,sec:empirical}). We state its limits with equal emphasis: our eight datasets are bimodal in $\lps$, the threshold is stress-tested in effectively one case, and within the exogenous regime $\lps$ does not order the size of the effect. Code, curation pipelines, logged results, and the pre-registration evidence trail are provided in an anonymized repository for review.\footnote{Anonymized for double-blind review; the reviewer link is provided separately to the editor. A de-anonymized public repository accompanies the camera-ready version.}
\end{itemize}

A word on framing. We do not claim that RevIN was ever asserted to be universally beneficial, and this paper is not an attack on it. The distribution-shift motivation of \citet{kim2022revin} is sound, and our results confirm that instance-normalization variants remain the best choice on the standard LTSF benchmarks. What we contribute is a \emph{characterization of the applicability boundary}: instance normalization rests on the assumption that lookback-window statistics extrapolate to the forecast horizon, and we state---in closed form, and via a pre-training diagnostic---when that assumption holds and when it fails.

The empirical picture is a double dissociation, not a ranking. The load-bearing comparison is the information-fair one: when every arm receives the identical past and future covariates, instance normalization is worse in mean MSE than even global $z$-scoring (\rawm) on all five covariate-capable backbones, so what separates the arms is the layer and not what the layer is allowed to see. The same dissociation appears at the level of strategies rather than of one method: against seasonal-naive and climatology anchors, covariate-conditioned level models---a first-stage regression alone, a classical dynamic regression, and \cnm---beat the best instance-normalized arm in 11 of 12 comparisons on the exogenously driven series and in 0 of 12 on the standard benchmarks, and the record is unchanged under MASE. Variance attribution locates the effect precisely: the normalization$\times$dataset interaction explains $47.1\%$ of log-MSE variance against $0.6\%$ for all backbone-related terms. That share is carried by the exogenous arm, and this is the informative fact rather than an artifact---restricted to the four \emph{endogenous} arms the same interaction is $0.4\%$ and all normalization-related terms total $0.7\%$, so the endogenous family is internally near-homogeneous and the boundary runs between endogenous and exogenous level handling, not within the endogenous family. The pre-training $\lps$ separates the two regimes: its sign prediction is correct on all eight datasets under MSE, MAE and MASE alike. We are deliberate about the limits of that separation---the eight $\lps$ values are bimodal, only one dataset lies near the threshold, and within the exogenous regime $\lps$ does not order the size of the gap (\Cref{sec:empirical}).

The paper proceeds as follows. \Cref{sec:related} positions our work in the normalization and benchmarking literatures. \Cref{sec:theory} develops the theory (C1); \Cref{sec:lps} defines the $\lps$ diagnostic and the pre-registration procedure (C3). \Cref{sec:synthetic,sec:empirical} present the synthetic and real-data validation (C2). \Cref{sec:discussion} discusses scope and limitations, and \Cref{sec:conclusion} concludes. Proofs are in \Cref{app:proofs}; design principles, reproducibility details, and extended results are in \Cref{app:design,app:repro,app:ext}.
